% 数学ノート集・共通テンプレート v2（積分論2026年度第1部の版面を参考）
% A5判。章内で定義・定理・例・注意等の番号を共通化。
% コンパイル： platex -interaction=nonstopmode math_note_template_v2.tex (2回)
%              dvipdfmx math_note_template_v2.dvi
\documentclass[a5paper,10pt,onecolumn]{jsarticle}
\usepackage[dvipdfmx,top=22mm,bottom=20mm,inner=19mm,outer=18mm,headheight=13pt,headsep=6mm]{geometry}
\usepackage{amsmath,amssymb,amsthm,mathtools}
\usepackage{enumitem,fancyhdr,titlesec}
\usepackage[dvipdfmx,bookmarks=true,colorlinks=false]{hyperref}
\usepackage{pxjahyper}

% ----- 差し替える文書情報 -----
\newcommand{\NoteTitle}{局所コンパクト性とパラコンパクト性}
\newcommand{\NoteSeries}{数学ノート集}

% ----- 文字・組版 -----
% 参考PDFの質素な明朝＋Computer Modernの印象を引き継ぎながら，
% A5判で窮屈にならないよう本文の行間を少し確保する．
\linespread{1.13}
\setlength{\parindent}{1zw}
\setlength{\parskip}{0pt}
\setlength{\abovedisplayskip}{6pt plus 2pt minus 2pt}
\setlength{\belowdisplayskip}{6pt plus 2pt minus 2pt}
\setlength{\abovedisplayshortskip}{3pt plus 2pt minus 1pt}
\setlength{\belowdisplayshortskip}{4pt plus 2pt minus 1pt}
\setlist[enumerate]{topsep=3pt,itemsep=1pt,parsep=0pt,leftmargin=2.4zw}
\titleformat{\section}{\normalfont\normalsize\bfseries}{\thesection}{0.8zw}{}
\titlespacing*{\section}{0pt}{2.0ex plus .4ex}{.9ex}
\titleformat{\subsection}{\normalfont\normalsize\bfseries}{\thesubsection}{0.7zw}{}
\titlespacing*{\subsection}{0pt}{1.5ex plus .3ex}{.6ex}
\raggedbottom

% ----- 柱とページ番号：ページ上部、罫線なし -----
\fancypagestyle{mathnote}{%
  \fancyhf{}%
  \fancyhead[L]{\footnotesize \NoteSeries\ （\NoteTitle）}%
  \fancyhead[R]{\footnotesize \thepage}%
  \renewcommand{\headrulewidth}{0pt}%
  \renewcommand{\footrulewidth}{0pt}%
}
\pagestyle{mathnote}

% ----- 節内通し番号の定義・定理群 -----
\newtheoremstyle{mathnotetheorem}
 {4.5pt}{4.5pt}{\normalfont}{}%
 {\bfseries}{}{0.65zw}%
 {\thmname{#1}\thmnumber{ #2}\thmnote{（#3）}}
\theoremstyle{mathnotetheorem}
\newtheorem{theorem}{定理}[section]
\newtheorem{definition}[theorem]{定義}
\newtheorem{lemma}[theorem]{補題}
\newtheorem{proposition}[theorem]{命題}
\newtheorem{corollary}[theorem]{系}
\newtheorem{example}[theorem]{例}
\newtheorem{remark}[theorem]{注意}
\renewcommand{\proofname}{証明}
\renewcommand{\qedsymbol}{$\square$}
\makeatletter
\renewenvironment{proof}[1][\proofname]{%
  \par\pushQED{\qed}\normalfont\topsep4\p@\@plus2\p@\relax
  \trivlist\item[\hskip\labelsep\bfseries #1]\ignorespaces
}{\popQED\endtrivlist\@endpefalse}
\makeatother

% ----- 使うノートで適宜拡張する記号 -----
\newcommand{\R}{\mathbb{R}}
\newcommand{\norm}[1]{\lVert #1\rVert}
\newcommand{\abs}[1]{\lvert #1\rvert}

\begin{document}
\section*{読了後にできること}
本稿を読むための前提は，局所コンパクト性，開被覆，単位の分割．
本稿の具体的な到達課題は，具体的な空間が局所コンパクトか判定し，局所有限な構造を使う．



\par\noindent\textbf{記号と約束.}\quad
関数空間の添字 $c$ はコンパクトな台をもつことを表す．
\par

\section{局所コンパクト性と正則性}
\begin{definition}
Hausdorff空間 $X$ が局所コンパクトであるとは，
各点がコンパクトな閉包を持つ近傍を持つことをいう．
開被覆 $\mathcal U$ の開細分とは，
各元が $\mathcal U$ のある元に含まれる開被覆をいう．
開細分が局所有限であるとは，
各点がそのうち有限個としか交わらない近傍を持つことをいう．
全ての開被覆が局所有限な開細分を持つ空間を
パラコンパクトという．
\end{definition}
\begin{lemma}[相対コンパクトな開近傍]\label{lem:relcompact}
$X$ が局所コンパクトHausdorffで
$x\in U$（$U$ 開）なら
$x\in V\subset\overline V\subset U$ で
$\overline V$ がコンパクトとなる開集合 $V$ がある．
\end{lemma}
\begin{proof}
$x$ の相対コンパクトな開近傍 $W$ を取る．
$\overline W$ はコンパクトHausdorffなので正規であり，
$x\in W\cap U$ の開近傍 $V$ を
$\overline V\subset W\cap U$ となるように
$\overline W$ 内で選べる．
$V$ を $X$ の開集合へ必要に応じ縮めれば
$\overline V$ は $\overline W$ の閉集合でありコンパクトである．
\end{proof}
\section{コンパクトな尽くし}
\begin{theorem}[可算コンパクト尽くし]\label{thm:exhaust}
局所コンパクトHausdorff第二可算空間 $X$ には，
コンパクト集合列 $(K_n)$ で
\[
 K_n\subset\operatorname{int}K_{n+1},\qquad
 X=\bigcup_{n=1}^\infty\operatorname{int}K_n
\]
を満たすものが存在する．
\end{theorem}
\begin{proof}
補題\ref{lem:relcompact}によって
相対コンパクトな開集合は $X$ の基底をなす．
第二可算性からそのうち可算個
$W_1,W_2,\dots$ で $X$ を覆うものを選べる．
帰納的に
$\overline{W_1}\cup\cdots\cup\overline{W_n}
 \cup K_{n-1}$
を内部に含む相対コンパクトな開集合の有限和を $V_n$ とする．
これはコンパクト集合の各点について補題を適用し，
有限部分被覆を取れば可能である．
$K_n=\overline{V_n}$ と置けばコンパクトで，
$K_{n-1}\subset V_n=\operatorname{int}K_n$ を満たし，
$\bigcup\operatorname{int}K_n$ は全ての $W_n$ を含む．
\end{proof}
\section{パラコンパクト性}
\begin{theorem}[局所コンパクト第二可算空間のパラコンパクト性]
\label{thm:para}
局所コンパクトHausdorff第二可算空間はパラコンパクトである．
\end{theorem}
\begin{proof}
定理\ref{thm:exhaust}のコンパクト尽くしを取り，
$K_0=K_{-1}=\varnothing$ とする．
\[
 A_n=K_n\setminus\operatorname{int}K_{n-1},
 \qquad W_n=\operatorname{int}K_{n+1}\setminus K_{n-2}
\]
と置く．$A_n$ はコンパクトで $A_n\subset W_n$ である．
与えられた開被覆 $\mathcal U$ に対し，
各 $A_n$ の点の周りで，
$\mathcal U$ のある元と $W_n$ の双方に閉包まで含まれる
相対コンパクト開近傍を取り，有限個で $A_n$ を覆う．
それらの全体を $\mathcal V$ とする．
各 $A_n$ が覆われ，$\bigcup_nA_n=X$ なので
$\mathcal V$ は開細分である．
$W_n$ は局所有限である：
$x\in\operatorname{int}K_N$ なら，
$n\ge N+2$ では $x\in K_{n-2}$ となり
$x\notin W_n$ である．
より小さい $\operatorname{int}K_N$ を近傍として取れば
$W_n$ の十分大きい添字の元とは交わらない．
各 $W_n$ に属する選択近傍は有限個なので
$\mathcal V$ も局所有限となる．
\end{proof}
\section{単位の分割}
\begin{lemma}[コンパクト集合の切断関数]
局所コンパクトHausdorff空間のコンパクト集合 $K\subset U$
（$U$ 開）に対し，
$f\in C_c(X,[0,1])$ で $f=1$ on $K$，
$\operatorname{supp}f\subset U$ となるものがある．
\end{lemma}
\begin{proof}
各 $x\in K$ について補題\ref{lem:relcompact}により
$x\in V_x\subset\overline{V_x}\subset U$ で
$\overline{V_x}$ がコンパクトとなるものを取る．
有限個の $V_x$ が $K$ を覆う．
それらの閉包の和を $L$ とすれば $L$ は
コンパクトHausdorffゆえ正規である．
$K$ と $L\setminus\bigcup V_x$ は互いに素な閉集合であり，
Urysohnの補題によって $L$ 上の連続関数で
$K$ 上1，$L\setminus\bigcup V_x$ 上0となるものが存在する．
この関数を $L$ の外で0に延長すれば連続で，
台は $U$ のコンパクト部分集合に含まれる．
\end{proof}
\begin{theorem}[連続な単位の分割]\label{thm:partition}
$X$ を局所コンパクトHausdorff第二可算空間とする．
任意の開被覆 $\mathcal U$ に対して，
局所有限な非負連続関数族 $(\varphi_i)$ で
\[
 \sum_i\varphi_i=1,\qquad
 \operatorname{supp}\varphi_i\subset U_i
 \quad(U_i\in\mathcal U)
\]
を満たすものが存在する．
\end{theorem}
\begin{proof}
定理\ref{thm:para}の証明で構成したコンパクト集合
$A_n$ と局所有限開集合 $W_n$ を用いる．
各 $A_n$ の点に対し，$\mathcal U$ のある元に含まれ
かつ $W_n$ に閉包まで含まれる相対コンパクト開近傍
$V_{n,j}$ を選び，有限個で $A_n$ を覆う．
各 $V_{n,j}$ の中に閉包まで含まれる開集合 $O_{n,j}$ を
補題\ref{lem:relcompact}で選び，
有限個の $O_{n,j}$ でも $A_n$ を覆うようにしてよい．
$K_{n,j}=\overline{O_{n,j}}$ は
$V_{n,j}$ に含まれるコンパクト集合である．
前補題を $K_{n,j}\subset V_{n,j}$ に適用して
$f_{n,j}\in C_c(X,[0,1])$ を選び，
$f_{n,j}=1$ on $K_{n,j}$，
$\operatorname{supp}f_{n,j}\subset V_{n,j}$ とする．
各 $n$ に対して関数は有限個であり，
$V_{n,j}\subset W_n$，$(W_n)$ は局所有限なので，
関数族 $(f_{n,j})$ も局所有限である．
各 $A_n$ は $(O_{n,j})_j$ で覆われ，
$X=\bigcup_nA_n$ なので全ての点で少なくとも一つの
$f_{n,j}$ が1となる．
局所有限性から $F=\sum_{n,j} f_{n,j}$ は連続であり，
全点で少なくとも一つの $f_{n,j}=1$ なので $F\ge1$．
$\varphi_{n,j}=f_{n,j}/F$ と置くと連続・非負であり，
局所有限性，台の包含および和が1であることが成立する．
\end{proof}
\begin{remark}
滑らかな多様体上の滑らかな単位の分割は，
局所的に滑らかなbump関数を用いて同じ正規化を行うことで得られる．
\end{remark}

\section{到達課題の完全解答}
\noindent\textbf{局所コンパクト性の具体的な判定．}
\begin{proposition}[Euclid空間の開集合]
$U\subset\mathbb R^n$ を開集合とすると，
$U$ は局所コンパクトである．
\end{proposition}
\begin{proof}
任意の $x\in U$ に対して
$\overline{B(x,r)}\subset U$ となる $r>0$ を取れる．
$\overline{B(x,r/2)}$ は $\mathbb R^n$ の有界閉集合なので
Heine--Borelの定理によりコンパクトである．
$B(x,r/2)$ はそのコンパクトな閉包を持つ開近傍である．
\end{proof}
\begin{proposition}[有理数の非局所コンパクト性]
$\mathbb Q$ に $\mathbb R$ からの相対位相を入れると
$\mathbb Q$ は局所コンパクトでない．
\end{proposition}
\begin{proof}
任意の $q\in\mathbb Q$ の近傍には
$J=(q-\varepsilon,q+\varepsilon)\cap\mathbb Q$
を含むものを取れる．
もし $q$ にコンパクトな近傍 $K\subset\mathbb Q$
があれば，ある $\varepsilon>0$ で $J\subset K$．
包含 $\mathbb Q\hookrightarrow\mathbb R$ は連続なので
$K$ は $\mathbb R$ でコンパクト，したがって閉集合である．
すると $\overline J^{\,\mathbb R}$ が $K$ に含まれるが，
この閉包には無理数が含まれ矛盾する．
\end{proof}
\noindent\textbf{到達確認．}
$\mathbb R^n$ の閉部分集合は局所コンパクトか．
\par\noindent\textbf{解答．}
閉集合 $F$ の任意の点 $x$ について
$F\cap\overline{B(x,1)}$ はコンパクトであり，
$F\cap B(x,1)$ は $F$ の開近傍である．
よって局所コンパクトである．
\end{document}
